\chapter{La protection des écosystèmes : étude d'une forêt}

\noindent\textit{Une forêt n'est pas une simple collection d'arbres : c'est un
ensemble d'êtres vivants en interaction permanente avec leur milieu. Ce dernier
chapitre étudie, à travers l'exemple d'un écosystème forestier, la biodiversité,
l'interdépendance entre espèces, les menaces d'origine humaine, et les moyens de
restaurer et protéger ces milieux.}
\vspace{8pt}

\connecteBanner

\section{Qu'est-ce qu'un écosystème forestier ?}

\begin{definition}[Écosystème]
Un \textbf{écosystème} est l'ensemble formé par un \textbf{milieu} (le sol, le
climat, l'eau disponible\dots) et tous les \textbf{êtres vivants} qui y habitent, en
interaction les uns avec les autres et avec leur milieu.
\end{definition}

\begin{propriete}[Une forêt, un écosystème riche]
Une forêt rassemble une grande variété d'êtres vivants : arbres et plantes du
sous-bois, insectes, oiseaux, mammifères, champignons, et une multitude de
microorganismes du sol. Chacun occupe une place précise dans l'écosystème.
\end{propriete}

\section{Biodiversité et interdépendance dans la forêt}

\begin{definition}[Biodiversité]
La \textbf{biodiversité} d'un milieu est la diversité des êtres vivants qui le
peuplent : diversité des espèces végétales, animales, fongiques et microbiennes.
\end{definition}

\begin{propriete}[Un réseau d'interdépendances]
Les espèces d'une forêt \textbf{dépendent les unes des autres} :
\begin{itemize}
\item les plantes servent de nourriture aux herbivores, eux-mêmes proies de
prédateurs (\textbf{chaîne alimentaire}) ;
\item de nombreux insectes assurent la \textbf{pollinisation} des fleurs ;
\item certains animaux dispersent les \textbf{graines}, participant au
renouvellement de la forêt ;
\item les \textbf{décomposeurs} (champignons, bactéries du sol) recyclent la matière
organique morte (feuilles, bois, cadavres) en éléments réutilisables par les
plantes.
\end{itemize}
\end{propriete}

\begin{illus}
\begin{tikzpicture}[scale=0.85]
\foreach \x/\c/\lbl in {0/onbgreen/{Feuilles}, 3.1/onbo/{Chenille},
  6.2/onbo2/{Oiseau}, 9.3/onbink/{Prédateur}}{
  \draw[\c,fill=\c,fill opacity=0.15,line width=1.2pt] (\x,0) circle (0.85);
  \node[font=\scriptsize,align=center] at (\x,0) {\lbl};
}
\draw[->,onbmuted,line width=1pt] (0.95,0) -- (2.25,0);
\draw[->,onbmuted,line width=1pt] (4.05,0) -- (5.35,0);
\draw[->,onbmuted,line width=1pt] (7.15,0) -- (8.45,0);
\end{tikzpicture}
\fig{Figure — Une chaîne alimentaire simple au sein d'un écosystème forestier.}
\end{illus}

\begin{exemple}
Si les décomposeurs venaient à disparaître, la matière organique morte
s'accumulerait sans être recyclée, privant les plantes des éléments nutritifs
nécessaires à leur croissance.
\end{exemple}

\begin{remarque}
Piège fréquent : la disparition d'une seule espèce peut déséquilibrer tout
l'écosystème, car de nombreuses espèces en dépendent directement ou indirectement
(pour se nourrir, se reproduire ou se protéger).
\end{remarque}

\section{Les activités humaines qui menacent la forêt}

\begin{propriete}[Principales menaces]
\begin{itemize}
\item la \textbf{déforestation} : exploitation excessive du bois, défrichement pour
l'agriculture ou l'urbanisation ;
\item le \textbf{braconnage} et la chasse excessive de certaines espèces animales ;
\item la \textbf{pollution} des sols et des cours d'eau ;
\item l'introduction d'\textbf{espèces invasives}, qui concurrencent les espèces
locales.
\end{itemize}
\end{propriete}

\begin{propriete}[Conséquences]
Ces activités entraînent une \textbf{perte de biodiversité}, la rupture de chaînes
alimentaires entières, une érosion accrue des sols (les racines des arbres ne
retiennent plus la terre), et des perturbations climatiques locales.
\end{propriete}

\begin{exemple}
La déforestation d'une zone entraîne la disparition des arbres qui servaient
d'habitat et de nourriture à de nombreuses espèces animales, provoquant leur
migration ou leur disparition locale.
\end{exemple}

\section{Restaurer et protéger les écosystèmes forestiers}

\begin{propriete}[Des mesures de protection]
\begin{itemize}
\item la création de \textbf{parcs nationaux} et de \textbf{réserves naturelles}, où
l'exploitation est interdite ou strictement encadrée ;
\item des \textbf{lois} protégeant les espèces menacées et limitant la chasse et
l'exploitation forestière ;
\item le \textbf{reboisement} des zones dégradées ;
\item les \textbf{jardins botaniques et zoologiques}, qui conservent et étudient des
espèces végétales et animales, parfois rares, et participent à la sensibilisation du
public.
\end{itemize}
\end{propriete}

\begin{exemple}
La réserve de faune du Dja, au Cameroun, est un vaste massif de forêt tropicale
protégée, qui abrite une grande diversité d'espèces végétales et animales.
\end{exemple}

\begin{remarque}
La protection d'un écosystème forestier ne se limite pas à interdire toute activité
humaine : elle vise à concilier la préservation de la biodiversité avec les besoins
des populations qui vivent de la forêt, par une gestion durable des ressources.
\end{remarque}

% ═══════════════════════════════════════════════════════════════
\section{L'essentiel à retenir}

\begin{synthese}
\textbf{Écosystème.} Milieu + êtres vivants en interaction.\\[4pt]
\textbf{Biodiversité forestière.} Diversité des espèces végétales, animales,
fongiques et microbiennes d'une forêt.\\[4pt]
\textbf{Interdépendance.} Chaînes alimentaires, pollinisation, dispersion des
graines, décomposeurs recyclant la matière organique.\\[4pt]
\textbf{Menaces humaines.} Déforestation, braconnage, pollution, espèces
invasives.\\[4pt]
\textbf{Protection.} Parcs nationaux, réserves naturelles, lois, reboisement,
jardins botaniques et zoologiques.\\[4pt]
\textbf{Piège fréquent.} La disparition d'une seule espèce peut déséquilibrer tout
l'écosystème, du fait des nombreuses interdépendances entre espèces.
\end{synthese}

% ═══════════════════════════════════════════════════════════════
\section{Méthodes \& savoir-faire}

\methode{Reconnaître une relation d'interdépendance}
Identifier comment une espèce dépend d'une autre pour se nourrir, se reproduire ou se
protéger.
\begin{exemple}
Un insecte pollinisateur et la plante qu'il féconde sont en relation
d'interdépendance.
\end{exemple}

\methode{Relier une activité humaine à une conséquence sur l'écosystème}
Identifier le mécanisme en jeu : destruction d'habitat (déforestation), prélèvement
excessif (braconnage), rejets nocifs (pollution).
\begin{exemple}
Le défrichement d'une forêt pour l'agriculture est une cause de déforestation,
entraînant la disparition d'habitats.
\end{exemple}

% ═══════════════════════════════════════════════════════════════
\section{Exercices d'application}
\begin{multicols}{2}

\rubrique{Écosystème forestier}
\exo{}\diff{1} Définir un écosystème.

\exo{}\diff{2} Citer trois catégories d'êtres vivants présents dans une forêt.

\rubrique{Interdépendance}
\exo{}\diff{2} Qu'est-ce qu'une chaîne alimentaire ?

\exo{}\diff{2} Quel est le rôle des décomposeurs dans un écosystème forestier ?

\rubrique{Menaces humaines}
\exo{}\diff{1} Citer deux menaces humaines qui pèsent sur les forêts.

\exo{}\diff{2} Quelle conséquence la déforestation a-t-elle sur les animaux qui
vivaient dans la forêt ?

\rubrique{Protection}
\exo{}\diff{1} Citer deux moyens de protéger un écosystème forestier.

\exo{}\diff{2} Quel est le rôle d'un jardin botanique ou zoologique ?

% ═══════════════════════════════════════════════════════════════
\end{multicols}

\section{Exercices d'entraînement}
\begin{multicols}{2}

\rubrique{Interdépendance}
\exo{}\diff{3} Expliquer pourquoi la disparition des insectes pollinisateurs
menacerait de nombreuses espèces végétales.

\exo{}\diff{3} Expliquer ce qui se passerait dans une forêt si tous les décomposeurs
disparaissaient.

\rubrique{Menaces humaines}
\exo{}\diff{3} Expliquer pourquoi la déforestation favorise l'érosion des sols.

\exo{}\diff{2} Pourquoi une espèce invasive peut-elle déséquilibrer un écosystème
forestier ?

\rubrique{Protection}
\exo{}\diff{3} Expliquer pourquoi la protection d'une forêt doit aussi prendre en
compte les besoins des populations qui en vivent.

\exo{}\diff{2} En quoi une réserve naturelle diffère-t-elle d'un jardin zoologique
dans sa façon de protéger les espèces ?

% ═══════════════════════════════════════════════════════════════
\end{multicols}

\section{Sujets type BEPC}

\sujetbac[Biodiversité et interdépendance forestière]
\begin{enumerate}[label=\arabic*)]
\item Définis un écosystème.
\item Cite trois exemples d'interdépendance entre espèces dans une forêt.
\item Explique le rôle des décomposeurs.
\item Explique pourquoi la disparition d'une espèce peut avoir des conséquences sur
plusieurs autres espèces.
\end{enumerate}

\sujetbac[Menaces et protection des forêts]
\begin{enumerate}[label=\arabic*)]
\item Cite trois menaces humaines qui pèsent sur les écosystèmes forestiers.
\item Cite deux conséquences de la déforestation.
\item Cite trois moyens de protéger ou de restaurer un écosystème forestier.
\item Pourquoi la gestion d'une forêt protégée doit-elle concilier préservation de
la biodiversité et besoins des populations locales ?
\end{enumerate}

% ═══════════════════════════════════════════════════════════════
\section{Corrigés}
\begin{multicols}{2}

\corrige{de l'exercice 1}
L'ensemble formé par un milieu et les êtres vivants qui y habitent, en interaction.

\corrige{de l'exercice 2}
Par exemple : plantes, insectes, champignons (ou mammifères, oiseaux).

\corrige{de l'exercice 3}
Une suite d'êtres vivants où chacun se nourrit du précédent (plante $\to$ herbivore
$\to$ prédateur).

\corrige{de l'exercice 4}
Ils recyclent la matière organique morte en éléments réutilisables par les plantes.

\corrige{de l'exercice 5}
Par exemple : la déforestation et le braconnage.

\corrige{de l'exercice 6}
Ils perdent leur habitat et leur source de nourriture, ce qui provoque leur migration
ou leur disparition locale.

\corrige{de l'exercice 7}
Par exemple : la création d'une réserve naturelle et le reboisement.

\corrige{de l'exercice 8}
Conserver et étudier des espèces, parfois rares, et sensibiliser le public à leur
protection.

\corrige{de l'exercice 9}
Car de nombreuses plantes dépendent des insectes pour être fécondées ; sans
pollinisateurs, elles ne pourraient plus se reproduire normalement.

\corrige{de l'exercice 10}
La matière organique morte (feuilles, bois, cadavres) s'accumulerait sans être
recyclée, privant les plantes des éléments nutritifs nécessaires à leur croissance.

\corrige{de l'exercice 11}
Car les racines des arbres, qui retiennent normalement la terre, disparaissent avec
les arbres abattus, rendant le sol plus vulnérable au ruissellement et au vent.

\corrige{de l'exercice 12}
Car elle concurrence les espèces locales pour les ressources (nourriture, espace),
sans prédateur naturel pour limiter sa prolifération.

\corrige{de l'exercice 13}
Car de nombreuses populations tirent leurs ressources (bois, nourriture, revenus) de
la forêt ; une protection trop stricte, sans alternative, pourrait leur nuire.

\corrige{de l'exercice 14}
Une réserve naturelle protège les espèces dans leur milieu naturel ; un jardin
zoologique les conserve et les étudie en dehors de leur milieu d'origine.

\corrige{du sujet 1}
1) L'ensemble formé par un milieu et les êtres vivants qui y habitent, en
interaction.\\
2) Par exemple : pollinisation, dispersion des graines, chaîne alimentaire.\\
3) Ils recyclent la matière organique morte, la rendant à nouveau disponible pour les
plantes.\\
4) Car de nombreuses espèces dépendent directement ou indirectement de l'espèce
disparue pour se nourrir, se reproduire ou se protéger.

\corrige{du sujet 2}
1) Par exemple : déforestation, braconnage, pollution.\\
2) Par exemple : perte de biodiversité et érosion des sols.\\
3) Par exemple : réserves naturelles, lois de protection, reboisement.\\
4) Car de nombreuses populations locales dépendent des ressources de la forêt pour
vivre ; une gestion durable permet de préserver la forêt tout en répondant à ces
besoins.

\end{multicols}
\exercicesBanner
